\begin{circuitikz}
    \draw (0,0) rectangle node[above=40pt]{\underline{Faktaboks}} (4,4);
    \draw ++(.5,2.5) node[above]{$R_x$} to[R, v_=$V_x$] +(3,0) node[above left]{$\rightarrow I_x$};
    \draw[<-] ++(1.5,3) .. controls +(.5,-.1) .. +(1,0) node[above]{$I_B$};
    \draw[->] ++(1.5,1.5) .. controls +(.5,.1) .. +(1,0) node[right]{$I_A$};
    \draw (.3,1) node[right]{$I_X=I_A-I_B$};
    \draw (.2,.5) node[right]{$V_X=R_X(I_A-I_B)$};
\end{circuitikz}
\begin{enumerate}
    \item Tegn en sløjfe, \slojfe[->], i alle masker, og navngiv maskestrømmene, $I_A,I_B,\ldots$
    \item For alle \emph{strømkilder}, opskriv en ligning af typen:

	\begin{center}
	\begin{circuitikz}
	    \draw (0,0) rectangle node{$I_A-I_B=2A$} (2.5,.7);
	    \draw ++(0,2) to[isource] node[above]{$2\unit A$} +(2.5,0);
	    \draw[->] ++(.7,1.3) .. controls +(.25,.1) and +(-.25,.1) .. +(1,0) node[below]{$I_A$};
	    \draw[<-] ++(.7,2.7) .. controls +(.25,-.1) and +(-.25,-.1) .. +(1,0) node[above]{$I_B$};
	\end{circuitikz}
	\hspace{5em}
	\begin{circuitikz}
	    \draw (0,0) rectangle node{$I_A-I_B=3\unit{\frac AA}\cdot I_X$} (3.5,.7);
	    \draw ++(0,2) to[cisource] node[above]{$3\unit{\frac AA}\cdot I_X$} +(3.5,0);
	    \draw[->] ++(1.2,1.3) .. controls +(.25,.1) and +(-.25,.1) .. +(1,0) node[below]{$I_A$};
	    \draw[<-] ++(1.2,2.7) .. controls +(.25,-.1) and +(-.25,-.1) .. +(1,0) node[above]{$I_B$};
	\end{circuitikz}
	\end{center}
    \item For alle masker, der ikke indeholder en \emph{strømkilde}, opskriv en ligning af nedenstående type. Hvis en strømkilde sidder mellem to masker, laves \emph{en fælles} ligning uden om kilden.

	\hspace{-7em}\begin{circuitikz}
	    \draw (0,0) -- +(-.2,0) +(0,0) -- +(0,-.2) +(0,0)
		to[short, *-] ++(0,1) to[R=$R_1$] ++(0,1) to[short, -*] ++(0,1) -- +(-.2,0) +(0,0) -- +(0,.2) +(0,0)
		to[vsource, l_=$5V$] ++(3,0) to[cvsource, l_=$2\unit{\frac VV}\cdot V_X$, invert] ++(1,0) to[short, -*] ++(1,0) -- +(.2,0) +(0,0) -- +(0,.2) +(0,0)
		to[R=$R_2$] ++(0,-3) to[short, -*] ++(0,0) -- +(.2,0) +(0,0) -- +(0,-.2) +(0,0)
		to[R=$R_3$] ++(-5,0);
	    \draw[->,scale=.75] ++(2,1.5) .. controls +(0,0.555) and +(-0.555,0) .. ++(1,1)
		.. controls +(0.555,0) and +(0,0.555) .. ++(1,-1)
		.. controls +(0,-.555) and +(.555,0) .. ++(-1,-1) node[above]{$I_A$}; 
	    \draw[->] ++(1.7,-1) .. controls +(.75,.1) .. +(1.5,0) node[right]{$I_C$};
	    \draw[->] ++(6,.7) .. controls +(-.1,.75) .. +(0,1.5) node[right]{$I_B$};
	    \draw ++(6.5,1) rectangle node{$R_1(I_A-0)-5\unit{V}+2\unit{\frac VV}\cdot V_X+R_2(I_A-I_B)+R_3(I_A-I_C)=0$} +(11,.7);
	\end{circuitikz}
    \item For styrede kilder, opskriv et udtryk for styrestrømmen/styrespændingen (se faktaboksen)
    \item Løs ligningerne, $I_A=\ldots,I_B=\ldots,$ (brug gerne maple)
    \item Bestem de V,I eller P, du blev spurgt om (se faktaboksen) for spændingsforeskelle over strømkilder, benyttes KVL.
\end{enumerate}
